% ============================================================
% cita-cadenas-es.tex — CADENAS EN CASTELLANO COMUNES A TODOS LOS
% ESTILOS DE biblatex
% ------------------------------------------------------------
% SOLO LAS CADENAS QUE NI biblatex (spanish.lbx) NI EL ESTILO TRAEN:
% SON LA TRADUCCIÓN DE LOS CÓDIGOS DEL CATÁLOGO DE LA APLICACIÓN
% (TABLA cmb_biblatex: editor_type, thesis_type, pubstate) Y DE LOS
% AGREGADOS PROPIOS DE LA EDITORIAL. UNA CADENA ASÍ SOLO APARECE SI
% UNA ENTRADA LA USA: EN UN ESTILO QUE NO LA NECESITA NO CAMBIA NADA.
%
% LAS QUE spanish.lbx YA TRAE Y UN ESTILO REDEFINE (meses, url,
% phdthesis, compiler…) NO VAN ACÁ: SON DECISIONES DE CADA ESTILO Y
% QUEDAN EN SU cita-<estilo>-config.tex.
%
% LO CARGA QUIEN CARGA EL ESTILO, INMEDIATAMENTE DESPUÉS DE
% cita-<estilo>-config.tex (\NewBibliographyString EXIGE biblatex YA
% CARGADO):
%   libros  — preambulo-contrato.tex (PDF) Y EL DOCUMENTO MÍNIMO DE
%             preparar_citas_libro.sh (PASADA DE tex4ht), POR NOMBRE
%   revistas — m_XML, CON RUTA ABSOLUTA A ~/.gbpublisher/latex/
% NO LO CARGA EL CONFIG: \input RESUELVE CONTRA LA CARPETA DE LA
% COMPILACIÓN, NO CONTRA LA DEL ARCHIVO QUE LO PIDE, Y EN REVISTAS EL
% CONFIG SE CARGA POR RUTA ABSOLUTA DESDE OTRA CARPETA.
%
% \NewBibliographyString CON UNA CADENA QUE YA EXISTE NO HACE NADA
% (biblatex.sty); CON UNA NUEVA, LA DECLARA PARA TODOS LOS IDIOMAS.
% ============================================================

% SIN biblatex NO HAY CADENAS QUE DEFINIR: ES EL CASO DE vancouver,
% QUE VA POR bibtex (cita-vancouver-config.tex). EL ARCHIVO TERMINA
% ACÁ EN LUGAR DE CORTAR LA COMPILACIÓN CON \NewBibliographyString
% INDEFINIDO
\makeatletter
\@ifpackageloaded{biblatex}{\makeatother}{\makeatother\endinput}

% ------------------------------------------------------------
% 1. FUNCIONES DE EDITOR (editortype) QUE spanish.lbx NO TRAE
% ------------------------------------------------------------
% EL CATÁLOGO OFRECE coordinator Y direction. biblatex BUSCA LA
% CADENA DEL MISMO NOMBRE, SU PLURAL (…s) Y LA FORMA by…; SIN
% ELLAS IMPRIME EL CÓDIGO TAL CUAL («coordinator», MEDIDO EN IEEE).
% organizator ES UN ALIAS HEREDADO DE ISO 690: EL CATÁLOGO USA
% organizer, QUE spanish.lbx YA TRAE.
\NewBibliographyString{coordinator}
\NewBibliographyString{coordinators}
\NewBibliographyString{bycoordinator}
\NewBibliographyString{direction}
\NewBibliographyString{directions}
\NewBibliographyString{bydirection}
\NewBibliographyString{organizator}
\NewBibliographyString{organizators}
\NewBibliographyString{byorganizator}

% ------------------------------------------------------------
% 2. TIPOS DE TESIS Y TRABAJOS FINALES (thesis_type)
% ------------------------------------------------------------
% LOS CÓDIGOS DEL CATÁLOGO, SALVO mathesis, phdthesis Y candthesis,
% QUE spanish.lbx TRAE Y CADA ESTILO REDEFINE EN SU CONFIG. LOS
% ÚLTIMOS SEIS SON ALIAS HEREDADOS DE ISO 690, QUE EL CATÁLOGO NO
% OFRECE: QUEDAN PARA UN .bib ESCRITO A MANO.
\NewBibliographyString{bachelorthesis}
\NewBibliographyString{engineerthesis}
\NewBibliographyString{specializationthesis}
\NewBibliographyString{fieldwork}
\NewBibliographyString{finalessay}
\NewBibliographyString{finalmonograph}
\NewBibliographyString{finalproject}
\NewBibliographyString{gradothesis}
\NewBibliographyString{integrativefinalwork}
\NewBibliographyString{magisterthesis}
\NewBibliographyString{masterfinalwork}
\NewBibliographyString{postgradthesis}
\NewBibliographyString{teacherthesis}
\NewBibliographyString{researchproject}
\NewBibliographyString{residencyreport}
\NewBibliographyString{shortthesis}
\NewBibliographyString{technicalreport}
\NewBibliographyString{undergraduatethesis}
\NewBibliographyString{ingenieriathesis}
\NewBibliographyString{licenciaturathesis}
\NewBibliographyString{maestriathesis}
\NewBibliographyString{posgradothesis}
\NewBibliographyString{profesoradothesis}
\NewBibliographyString{tifthesis}

% ------------------------------------------------------------
% 3. ESTADOS DE PUBLICACIÓN (pubstate) Y FORMAS DE EDICIÓN
% ------------------------------------------------------------
\NewBibliographyString{mimeograph}
\NewBibliographyString{photocopy}
\NewBibliographyString{digitalprinting}
\NewBibliographyString{edauthor}
% self-published ES EL CÓDIGO DEL CATÁLOGO PARA LA EDICIÓN DE AUTOR (SIN
% CADENA, biblatex LO IMPRIMÍA TAL CUAL: MEDIDO EN VERONA E IEEE). edauthor
% QUEDA PARA UN .bib ESCRITO A MANO. EL GUION EN LA CLAVE ES VÁLIDO (MEDIDO)
\NewBibliographyString{self-published}
\NewBibliographyString{documentjob}
\NewBibliographyString{origpubbare}

% ------------------------------------------------------------
% 4. CINE: DIRECCIÓN, GUION Y ELENCO
% ------------------------------------------------------------
\NewBibliographyString{dirigida}
\NewBibliographyString{dirigidas}
\NewBibliographyString{bydirigida}
\NewBibliographyString{escrita}
\NewBibliographyString{escritas}
\NewBibliographyString{byescrita}
\NewBibliographyString{elenco}
\NewBibliographyString{elencos}
\NewBibliographyString{byelenco}

% ------------------------------------------------------------
% 5. «Y SIGUIENTE», «Y SIGUIENTES»: \sq Y \sqq EN UNA postnote
% ------------------------------------------------------------
\NewBibliographyString{sq}
\NewBibliographyString{sqq}
\providecommand*{\sq}{\bibstring{sq}}
\providecommand*{\sqq}{\bibstring{sqq}}

% ------------------------------------------------------------
% 6. EL TEXTO EN CASTELLANO
% ------------------------------------------------------------
\DefineBibliographyStrings{spanish}{%
  coordinator          = {coord\adddot},
  coordinators         = {coords\adddot},
  bycoordinator        = {coord\adddot\addspace por},
  direction            = {dir\adddot},
  directions           = {dirs\adddot},
  bydirection          = {dir\adddot\addspace por},
  organizator          = {org\adddot},
  organizators         = {orgs\adddot},
  byorganizator        = {org\adddot\addspace por},
  bachelorthesis       = {Tesis de Licenciatura},
  engineerthesis       = {Tesis de Ingeniería},
  specializationthesis = {Tesis de Especialización},
  fieldwork            = {Trabajo de Campo},
  finalessay           = {Ensayo Final},
  finalmonograph       = {Monografía Final de Carrera},
  finalproject         = {Proyecto Final},
  gradothesis          = {Tesis de Grado},
  integrativefinalwork = {Trabajo Integrador Final (TIF)},
  magisterthesis       = {Tesis de Magíster},
  masterfinalwork      = {Trabajo Final de Maestría},
  postgradthesis       = {Tesis de Posgrado},
  teacherthesis        = {Tesis de Profesorado},
  researchproject      = {Proyecto de Investigación},
  residencyreport      = {Informe Final de Residencia},
  shortthesis          = {Tesina},
  technicalreport      = {Trabajo Final Técnico},
  undergraduatethesis  = {Trabajo Final de Grado},
  ingenieriathesis     = {Tesis de Ingeniería},
  licenciaturathesis   = {Tesis de Licenciatura},
  maestriathesis       = {Tesis de Maestría},
  posgradothesis       = {Tesis de Posgrado},
  profesoradothesis    = {Tesis de Profesorado},
  tifthesis            = {Trabajo Integrador Final},
  mimeograph           = {mimeo},
  photocopy            = {fotocopia},
  digitalprinting      = {impresión digital},
  edauthor             = {edición de autor},
  self-published       = {edición de autor},
  documentjob          = {documento de trabajo},
  origpubbare          = {orig\adddotspace pub\adddotspace},
  dirigida             = {dirección},
  dirigidas            = {dirección},
  bydirigida           = {Dirección de},
  escrita              = {escrita},
  escritas             = {escrita},
  byescrita            = {escrita por},
  elenco               = {actúan:},
  elencos              = {actúan:},
  byelenco             = {actúan:},
  sq                   = {~y~sig\adddot},
  sqq                  = {~y~ss\adddot},
}
